\chapter{Testing and Continuous Delivery}\label{ch:pl:testing-and-continuous-delivery}

In the last days of July 2012 a broker-dealer deployed new order-routing code to its eight routing servers, in stages. The new code reused a
flag that had once switched on a function retired years before and never deleted. A technician did not copy the new code to one of the eight
servers; no second person reviewed the deployment, and no written procedure required one. On 1 August orders carrying the reused flag reached
the eighth server, woke the retired function, and in forty-five minutes the firm sent millions of orders and lost more than \$460 million. Book~11
reads the episode for its missing risk limits; this chapter reads it for its deployment. Every piece of it --- the untested old code, the
reused flag, the partial deployment, the missing check --- is something a test pipeline and a deployment procedure exist to catch.

\section{The test pyramid for trading code}

\begin{definition}[Unit, integration and end-to-end tests]\label{def:pl:testing-and-continuous-delivery:tests}
A \emph{unit test}\index{unit test} checks one function or class in isolation, in milliseconds. An \emph{integration test}\index{integration
test} checks several components working together --- a pricer with its market data, a strategy with the simulator --- in seconds. An
\emph{end-to-end test}\index{end-to-end test} runs a whole flow as production would --- orders from a strategy through the gateway to a
simulated venue and back to the \omterm{def:pl:position-and-pnl-services:service}{position service} --- in minutes.
\end{definition}

\begin{definition}[Test pyramid]\label{def:pl:testing-and-continuous-delivery:pyramid}
The \emph{test pyramid}\index{test pyramid} is the rule that a code base should have many \omterm{def:pl:testing-and-continuous-delivery:tests}{unit tests}, fewer \omterm{def:pl:testing-and-continuous-delivery:tests}{integration tests} and few
\omterm{def:pl:testing-and-continuous-delivery:tests}{end-to-end tests}, because each level up is slower, more fragile and less precise about what broke; it became widely known through Mike Cohn's
2009 book \emph{Succeeding with Agile}.
\end{definition}

Trading code adds levels of its own. Book~13, chapter~25, describes property-based, differential and fuzz tests, performance regression gates
and exchange certification for the low-latency path; chapter~\ref{ch:pl:simulation-production-parity}'s \omterm{def:pl:simulation-production-parity:harness}{parity harness} is an \omterm{def:pl:testing-and-continuous-delivery:tests}{integration test}
between the simulator and the production adapter; Book~6's model validation is a review, not a test, but it produces the reference values tests
use. This chapter adds the test a pricing library needs most, and the controls that decide what reaches production.

\section{Golden files for pricers}

\begin{definition}[Golden-file test]\label{def:pl:testing-and-continuous-delivery:golden}
A \emph{golden-file test}\index{golden-file test} compares a program's outputs with reviewed values stored with the code --- the golden
file --- each with a tolerance and the provenance of the value; any difference beyond the tolerance fails the test, and changing a golden value
is itself a reviewed change.
\end{definition}

A pricing library's outputs are numbers without an independent oracle for most products: a closed form covers the simplest, and Book~5's
reference pricers cover some more. What a golden file adds is memory: the value the library gave yesterday, reviewed then, against which
today's build is compared. The JPMorgan task force report on the 2012 losses in the Chief Investment Office found that the new VaR model ran
through spreadsheets filled by copying and pasting, and that one of them divided by the sum of two rates instead of their average, which
likely muted volatility by a factor of two. A golden value for the VaR of a reference portfolio, reviewed once, would have moved when the
model changed.

The chapter's portfolio has sixty instruments priced by Book~5's library: twenty European options (analytic engine), ten American puts
(finite differences), twenty Asian calls (Monte Carlo, 20\,000 paths) and ten down-and-out calls (closed form). The tolerance is the engine's
own error times a factor $k$: a relative $10^{-10}$ for the closed forms, $k$ times the gap between the 200-point and the 400-point grid for
the PDE, $k$ times the Monte Carlo standard error --- estimated from eight independent seeds --- for Monte Carlo.

\begin{omfigure}
\begin{tikzpicture}
\begin{semilogxaxis}[width=11cm,height=4.6cm,xbar,bar width=10pt,xmin=1e-10,xmax=30,xtick={1e-10,1e-8,1e-6,1e-4,1e-2,1},y dir=reverse,
  ytick=data,yticklabels={European (analytic),American (PDE),Asian (Monte Carlo),down-and-out (closed form)},
  xlabel={mean error unit per instrument (price units, log scale)},grid=major,grid style={gray!20},tick label style={font=\small},
  label style={font=\small},table/col sep=comma,log origin=infty,nodes near coords,point meta=rawx,
  every node near coord/.append style={font=\scriptsize,anchor=west,/pgf/number format/sci,/pgf/number format/precision=2}]
\addplot[fill=omDef!40,draw=omDef] table[y expr=\coordindex,x=mean_error_unit]
  {figdata/platforms/27-testing-and-continuous-delivery/errors.csv};
\end{semilogxaxis}
\end{tikzpicture}

\omcaption{The error unit behind each golden value, averaged by product: a relative $10^{-10}$ of prices near 12--15 for the closed forms, the
0.0019 gap between two PDE grids for American puts near 11, and the 0.053 Monte Carlo standard error of Asian calls near 8. The tolerance is
$k$ times this unit, so one $k$ means eight orders of magnitude between products. Data: \texttt{fig\_goldtest.py}.}
\label{fig:pl:testing-and-continuous-delivery:errors}
\end{omfigure}

\omcode{code/platforms/27-testing-and-continuous-delivery/python/pl_goldtest.py}{74}{82}{Each golden value's tolerance from its engine's
own error: Monte Carlo spread across seeds, grid convergence for the PDE, a relative epsilon for closed forms.}
\label{lst:pl:testing-and-continuous-delivery:tol}

\omcode{code/firm/goldtest/firm_goldtest.py}{80}{93}{The comparison: every result against its golden value and tolerance, with a missing value
on either side a failure, and the failures counted by product.}\label{lst:pl:testing-and-continuous-delivery:compare}

Five builds are then compared with the golden file. Each plants one change in the library's behaviour.

\omcode{code/platforms/27-testing-and-continuous-delivery/python/pl_goldtest.py}{103}{118}{The five planted builds: a new Monte Carlo seed, an
ACT/360 day count, a sign error in the rate of puts, an Asian average that drops its last fixing, and a finer PDE grid.}
\label{lst:pl:testing-and-continuous-delivery:builds}

\begin{omfigure}
\begin{tikzpicture}
\begin{axis}[width=11.5cm,height=6cm,xlabel={tolerance factor $k$ (Monte Carlo standard errors)},ylabel={share},
  grid=major,grid style={gray!20},tick label style={font=\small},label style={font=\small},table/col sep=comma,
  xmin=0.8,xmax=6.2,ymin=0,ymax=1.08,xtick={1,2,3,4,5,6}]
\addplot[omThm,thick,mark=*] table[x=k,y=asian_fix]{figdata/platforms/27-testing-and-continuous-delivery/tolerance.csv};
\addplot[black,thick,mark=square*] table[x=k,y=noise_run]{figdata/platforms/27-testing-and-continuous-delivery/tolerance.csv};
\addplot[omDef,thick,dashed,mark=triangle*] table[x=k,y=noise_instrument]{figdata/platforms/27-testing-and-continuous-delivery/tolerance.csv};
\addplot[gray,thin] table[x=k,y=noise_theory]{figdata/platforms/27-testing-and-continuous-delivery/tolerance.csv};
\node[omThm,font=\scriptsize,anchor=south west] at (axis cs:4.2,0.62) {Asian-fixing bug detected};
\node[font=\scriptsize,anchor=north west] at (axis cs:1.05,0.8) {$\blacksquare$ runs failed (noise)};
\node[omDef,font=\scriptsize,anchor=north west] at (axis cs:1.05,0.68) {$\blacktriangle$ Asian calls flagged};
\node[gray,font=\scriptsize,anchor=north east] at (axis cs:2.3,0.12) {theory};
\end{axis}
\end{tikzpicture}

\omcaption{The Monte Carlo tolerance trade-off. With the seed free to change, noise alone flags each Asian call with the probability that two
estimates differ by more than $k$ standard errors (theory in grey) and fails every run for $k\le 2$; the Asian-fixing bug, which only the Monte
Carlo instruments see, is caught on all twenty at $k\le2$ but on 9 of 20 at $k=5$, where noise still fails 2.5\% of runs. Data:
\texttt{fig\_goldtest.py}.}\label{fig:pl:testing-and-continuous-delivery:tolerance}
\end{omfigure}

The two large regressions are caught at every tolerance: the sign error changes all twenty puts by hundreds of error units and more, and the day
count changes every European, American and barrier price by far more than its tolerance. The trouble is Monte Carlo
(\cref{fig:pl:testing-and-continuous-delivery:tolerance}). With the seed free to change, the golden test faces two estimates of the same price
and must tolerate their difference. At $k=2$ noise alone flags 18\% of the Asian calls and fails every one of forty runs; at $k=5$ it fails 2.5\%
of runs, but the Asian-fixing bug --- a small change, two to ten standard errors --- is caught on only 9 of the 20 Asian calls, and the day
count's effect on them, never more than two standard errors, on none. With a false alarm costing \$1\,000 of investigation and a missed single-product
regression \$500\,000 --- the model's assumptions --- the cheapest tolerance is $k\le2$, a test that fails on every run. No tolerance works.

The fix is not a tolerance but a seed. Book~5's Monte Carlo engine draws its random numbers from a seed fixed per instrument; if the build
pins it, today's Monte Carlo price is a deterministic function of the code, the same random numbers give the same price to the last digit,
and the Monte Carlo instruments can be compared at a closed form's tolerance. Pinned, the golden test catches the Asian-fixing bug on all 20
instruments and the day count on all 60. A new seed then fails the test like any other change --- which it is --- and is released as an
approved update of the golden values. The same holds for the finer PDE grid: it moves each American put by exactly one error unit, passes
every $k\ge1$, and should reach production as a reviewed golden update with the new grid in its provenance, not as a silent pass.

\section{Continuous integration}

\begin{definition}[Continuous integration, deployment pipeline]\label{def:pl:testing-and-continuous-delivery:ci}
\emph{Continuous integration}\index{continuous integration} merges every change into the main line at least daily, and builds and tests the
merged code automatically on each merge. A \emph{deployment pipeline}\index{deployment pipeline} is the automated sequence of stages ---
build, \omterm{def:pl:testing-and-continuous-delivery:tests}{unit tests}, integration and golden tests, \omterm{def:pl:testing-and-continuous-delivery:tests}{end-to-end tests}, deployment --- that every change passes through, in order, stopping at the
first failure.
\end{definition}

\begin{definition}[Continuous delivery]\label{def:pl:testing-and-continuous-delivery:cd}
\emph{Continuous delivery}\index{continuous delivery} keeps the main line releasable at all times, so that any change that passes the
pipeline can be deployed to production by a decision rather than a project.
\end{definition}

\begin{omfigure}
\begin{tikzpicture}[font=\small,node distance=5mm,
  st/.style={draw=omDef,fill=omDef!8,rounded corners=2pt,align=center,minimum height=11mm,text width=19mm},
  gate/.style={draw=omThm,fill=omThm!8,rounded corners=2pt,align=center,minimum height=11mm,text width=19mm},
  arr/.style={-{Stealth[length=2mm]},thick}]
\node[st] (b) {build};
\node[st,right=of b] (u) {unit tests\\0.02~s};
\node[st,right=of u] (g) {golden tests\\0.33~s};
\node[st,right=of g] (e) {end to end\\1.62~s};
\node[gate,right=of e] (c) {every host\\consistent?};
\node[st,right=of c] (sw) {switch\\traffic};
\foreach \a/\b in {b/u,u/g,g/e,e/c,c/sw} \draw[arr] (\a) -- (\b);
\draw[arr,omThm] (c.south) -- ++(0,-6mm) -| node[pos=0.25,below,font=\scriptsize]{no: refuse, stay on blue} (b.south);
\end{tikzpicture}

\omcaption{The chapter's \omterm{def:pl:testing-and-continuous-delivery:ci}{deployment pipeline}, with the measured duration of each test tier on the sixty-instrument portfolio (otherwise idle machine;
\texttt{bench\_goldtest.py}). Each stage runs only if the one before passed; the last check before the traffic switch is the consistency of
every host with the release.}\label{fig:pl:testing-and-continuous-delivery:pipeline}
\end{omfigure}

The pipeline's order is the pyramid's: cheap tests first, so that most failures are found in seconds. On the chapter's portfolio, the unit tier
(put--call parity on a thousand options) takes 0.02~seconds, the golden tier (the sixty instruments against their pinned values) 0.33 and the
end-to-end tier (the portfolio revalued in five scenarios) 1.62, measured on an otherwise idle machine. The golden file is small enough to run on every
merge; a firm's real one, with thousands of instruments, runs in the nightly integration build and on every change to the library.

\section{Release trains and change control}

\begin{definition}[Release train]\label{def:pl:testing-and-continuous-delivery:train}
A \emph{release train}\index{release train} deploys, at fixed times, everything that has passed the pipeline since the previous departure; a
change that misses a train waits for the next.
\end{definition}

\begin{definition}[Change control]\label{def:pl:testing-and-continuous-delivery:change}
\emph{Change control}\index{change control} is the procedure every production change follows --- recorded, reviewed, approved by someone
other than its author, deployed by a stated method, verified --- so that the firm can say for any change who made it, who approved it and
when.
\end{definition}

The train's cadence sets a trade-off (\cref{tab:pl:testing-and-continuous-delivery:trains}). Over a model year of 533 changes, deploying each
change as it passes the pipeline gives no wait and one change per deployment; a daily train makes changes wait half a day on average and ships
1.8 at a time; a weekly train, 3.5 days and 10.3 changes; a fortnightly one, 7 days and 20.5. The batch matters when a deployment goes wrong:
finding the bad change among ten takes about four bisection steps, among twenty about five, and each step is a \omterm{def:pl:testing-and-continuous-delivery:bluegreen}{rollback}, a redeployment and a
wait for the symptom. Trading firms usually settle between the two: small changes often for research and analytics, trains with a change
window for the order path, where RTS~6 requires a clearly delineated methodology to develop and test a system before its deployment or
substantial update, in an environment separated from production.

\begin{table}[htbp]
\centering\small
\begin{tabular}{lrrr}
\toprule
cadence & mean wait (days) & changes per deployment & bisection steps\\
\midrule
each change at once & 0.00 & 1.00 & 0.00\\
daily train & 0.51 & 1.83 & 1.16\\
weekly train & 3.46 & 10.25 & 3.93\\
fortnightly train & 6.99 & 20.50 & 4.86\\
\bottomrule
\end{tabular}
\caption{A model year of 533 changes (ten a week) under four release cadences: the mean wait from passing the pipeline to production, the
changes per deployment, and the bisection steps needed, on average per change, to find a bad change in its deployment.}
\label{tab:pl:testing-and-continuous-delivery:trains}
\end{table}

\section{Deploying and rolling back}

\begin{definition}[Blue--green deployment, rollback]\label{def:pl:testing-and-continuous-delivery:bluegreen}
A \emph{blue--green deployment}\index{blue--green deployment} installs a release on a second, idle set of hosts (green) while the current set
(blue) serves, then switches traffic to green at once; switching back is the \emph{rollback}\index{rollback}, the return to the previous
release, which a blue--green deployment makes a switch rather than a reinstallation.
\end{definition}

The hook's deployment failed in the step between installation and switch: seven hosts ran the new release, one ran the old, and nothing
compared them. A deployment check is simple to write. Each host reports the hash of the build it runs and of its configuration, and the release
is refused --- the traffic switch does not happen --- unless every host reports the release's hashes. Run on the hook's eight servers, it names
the eighth before the market opens.

\omcode{code/firm/goldtest/firm_goldtest.py}{137}{145}{The deployment consistency check: each host's build and configuration hashes against
the release's; any mismatch refuses the release.}\label{lst:pl:testing-and-continuous-delivery:check}

\begin{omfigure}
\begin{tikzpicture}[font=\small,
  host/.style={draw=omDef,fill=omDef!8,rounded corners=2pt,minimum width=13mm,minimum height=9mm,align=center,font=\scriptsize},
  bad/.style={draw=omThm,fill=omThm!15,rounded corners=2pt,minimum width=13mm,minimum height=9mm,align=center,font=\scriptsize},
  arr/.style={-{Stealth[length=2mm]},thick}]
\foreach \i in {1,...,7} \node[host] (h\i) at (1.55*\i-1.55,0) {router\i\\8430046e};
\node[bad] (h8) at (1.55*7,0) {router8\\0290313a};
\node[draw=black,rounded corners=2pt,align=center,text width=44mm] (c) at (5.4,-2.1) {release 8430046e: check every host};
\foreach \i in {1,...,8} \draw[arr,gray] (h\i.south) -- (c.north);
\node[omThm,font=\scriptsize,align=left,anchor=west] at (8.3,-2.1) {refused: router8 runs\\another build};
\node[font=\scriptsize,anchor=south] at (5.4,0.6) {build hashes reported by the eight hosts (first 8 characters)};
\end{tikzpicture}

\omcaption{The consistency check on the hook's deployment, re-enacted: seven routers report the new build's hash, the eighth the old one, and
the release is refused before any order reaches the eighth. The same check on configuration hashes catches \omterm{def:pl:signal-serving-and-parameter-management:drift}{configuration drift}
(chapter~\ref{ch:pl:signal-serving-and-parameter-management}).}\label{fig:pl:testing-and-continuous-delivery:hosts}
\end{omfigure}

The order also lists what else was missing: the retired code was still present and callable, and the flag it had once used was given a new
meaning. Both are testing failures in this chapter's sense. Dead code should be deleted, not left behind a flag; a flag should never be
reused, since an old host or an old message will still give it the old meaning; and a \omterm{def:pl:signal-serving-and-parameter-management:flag}{feature flag}
(chapter~\ref{ch:pl:signal-serving-and-parameter-management}) switched on by a release is a deployment of its own, with its own staged
rollout and its own check that every host agrees on what it means.

\begin{dated}{2026-09}{The rules on testing and deployment}\label{dat:pl:testing-and-continuous-delivery:rules}
The SEC's order of 16 October 2013 in the matter of Knight Capital Americas LLC (Release 34-70694) found that new code repurposed a flag that
had activated a retired function still present on the servers; that a technician did not copy the new code to one of eight servers; that no
second technician reviewed the deployment and no written procedure required it; and that orders routed over forty-five minutes on 1 August
2012 cost the firm more than \$460 million. In the European Union, Commission Delegated Regulation (EU) 2017/589 (RTS~6) requires an
investment firm to establish clearly delineated methodologies to develop and test an algorithmic trading system before its deployment or
substantial update, ensuring among other things that it does not behave in an unintended manner, and to test in an environment separated
from production (articles 5 and 7).
\end{dated}

\section{Tutorial: seven of eight servers}\label{tut:pl:testing-and-continuous-delivery:tut}

\begin{tutorial}
\textbf{Goal.} Freeze a pricing library's golden values, measure what each tolerance detects and falsely flags, and refuse a partial
deployment. \textbf{End state:} \cref{fig:pl:testing-and-continuous-delivery:tolerance,tab:pl:testing-and-continuous-delivery:trains}.
\begin{tutsteps}
\item \textbf{Portfolio}: \texttt{pl\_goldtest.portfolio()}; the error unit of each instrument with \texttt{errors}.
\item \textbf{Builds}: \texttt{experiment()}, the five planted builds and forty noise builds.
\item \textbf{Tolerance}: \texttt{rates(ex, k)} for each $k$, then with \texttt{pinned=True}; \texttt{expected\_cost(ex, k)}.
\item \textbf{Pipeline}: \texttt{bench\_goldtest.py} for the stage times; \texttt{trains()} for the cadences.
\item \textbf{Deployment}: \texttt{eight\_servers()} and \texttt{firm\_goldtest.consistency}.
\end{tutsteps}
\textbf{What to change next.} Store the golden file on disk with \texttt{GoldenStore.save} and require an approval for the grid refinement;
add a check that no two releases give one flag two meanings.
\end{tutorial}

\section{Build: golden tests and deployment checks}\label{bld:pl:testing-and-continuous-delivery:goldtest}

\begin{build}
\bfield{Purpose} Know, for every build of the pricing library, which values moved and by how much, and never let a release reach a host set
that does not all run it.
\bfield{Interface} \texttt{Golden}, \texttt{GoldenStore} (\texttt{put}, \texttt{get}, \texttt{propose}, \texttt{approve}, \texttt{save},
\texttt{load}), \texttt{compare}, \texttt{run\_tiers}, \texttt{release\_train}, \texttt{host\_state}, \texttt{consistency}.
\bfield{Rules} Tolerances derived from engine error, not chosen; Monte Carlo seeds pinned in the build; golden updates proposed with a reason
and approved by someone else; tiers in pyramid order, stopping at the first failure; no release switched on until every host matches it.
\bfield{Acceptance tests} \texttt{code/firm/goldtest/tests/}: comparison with a failure, a missing value and a pass; approval by another
person and persistence; tiers stopping at a failure; release-train waits, batches and bisection; consistency with a stale build and a drifted
configuration.
\bfield{Stretch} A change-impact report by risk factor; golden values of risk (Greeks, VaR) as well as prices; blue--green switching in the
simulator.
\end{build}

\begin{omsources}
\item US Securities and Exchange Commission, \emph{In the Matter of Knight Capital Americas LLC}, Release No.~34-70694, 16 October 2013.
\item \emph{Report of JPMorgan Chase \& Co.\ Management Task Force Regarding 2012 CIO Losses}, 16 January 2013.
\item M.~Fowler, \emph{TestPyramid}, 2012 (on M.~Cohn, \emph{Succeeding with Agile}, 2009).
\item Commission Delegated Regulation (EU) 2017/589 (RTS~6), articles 5 and 7.
\item One Quant Book~5, chapter~28 (reference pricers); Book~11, chapter~27 (risk controls); Book~13, chapter~25 (testing low-latency code).
\end{omsources}

\section{Exercises}

\begin{exercise}[$\star$]\label{exo:pl:testing-and-continuous-delivery:1}
Why is the tolerance of a closed-form price a relative $10^{-10}$ and not zero?
\end{exercise}

\begin{exercise}[$\star$]\label{exo:pl:testing-and-continuous-delivery:2}
With the seed free, what share of Asian calls does noise alone flag at $k=3$, and what does theory predict?
\end{exercise}

\begin{exercise}[$\star$]\label{exo:pl:testing-and-continuous-delivery:3}
A weekly train carries ten changes and one is bad. About how many bisection steps does it take to find it?
\end{exercise}

\begin{exercise}[$\star\star$]\label{exo:pl:testing-and-continuous-delivery:4}
Why does the finer PDE grid pass the golden test for every $k\ge1$, and why should it still be approved?
\end{exercise}

\begin{exercise}[$\star\star$]\label{exo:pl:testing-and-continuous-delivery:5}
Why can the day count's effect on the Asian calls not be seen with the seed free, and how is it seen with the seed pinned?
\end{exercise}

\begin{exercise}[$\star\star$]\label{exo:pl:testing-and-continuous-delivery:6}
What would a \omterm{def:pl:testing-and-continuous-delivery:bluegreen}{blue--green deployment} have changed on the hook's morning, and what would it not have changed?
\end{exercise}

\begin{exercise}[$\star\star\star$]\label{exo:pl:testing-and-continuous-delivery:7}
\emph{Coding.} Double the Monte Carlo paths of the portfolio and rerun \texttt{rates(ex, 5)}. What happens to the Asian-fixing detection and
why?
\end{exercise}

\begin{exercise}[$\star\star\star$]\label{exo:pl:testing-and-continuous-delivery:8}
\emph{Find the flaw.} \textquotedbl{}The golden test failed again on Monte Carlo noise, so we widened every tolerance to ten standard
errors.\textquotedbl{}
\end{exercise}

\section{Problem: Seven of Eight Servers}

\begin{problem}[{Weekend problem --- seven of eight servers}]\label{pb:pl:testing-and-continuous-delivery:1}
The chapter's portfolio, builds, pipeline, trains and eight routers.

\medskip\noindent\textbf{Part I --- Tests.}
\begin{enumerate}
\item What does each level of the pyramid test, and why are there fewer tests at each level up?
\item What does a golden file hold for each value?
\item How is each engine's tolerance derived?
\item What did the JPMorgan task force report find in the VaR spreadsheets?
\item Which builds change which instruments?
\end{enumerate}

\medskip\noindent\textbf{Part II --- Tolerance.}
\begin{enumerate}[resume]
\item Which regressions are caught at every tolerance, and why?
\item What does noise alone do at $k=2$ and at $k=5$?
\item What share of the Asian-fixing bug is caught at each $k$?
\item Which $k$ minimises the model's expected cost, and what is wrong with it?
\item What does pinning the seed change?
\end{enumerate}

\medskip\noindent\textbf{Part III --- Delivery.}
\begin{enumerate}[resume]
\item How long does each tier of the pipeline take?
\item What do the four release cadences give?
\item What does \omterm{def:pl:testing-and-continuous-delivery:change}{change control} record?
\item What does RTS~6 require before a trading algorithm is deployed?
\item What is a \omterm{def:pl:testing-and-continuous-delivery:bluegreen}{rollback} in a \omterm{def:pl:testing-and-continuous-delivery:bluegreen}{blue--green deployment}?
\end{enumerate}

\medskip\noindent\textbf{Part IV --- The verdict.}
\begin{enumerate}[resume]
\item State the \emph{named result}: the detection rate of each planted regression and the false-alarm rate of Monte Carlo noise at each
tolerance, the tolerance that minimises expected cost, and the consistency check that would have refused the partial deployment.
\item What four things went wrong in the hook's deployment, in this chapter's terms?
\item Which one would the consistency check have caught, and which would testing have caught?
\item What would you require of every release of the order path?
\item In one sentence: what is a deployment?
\end{enumerate}
\end{problem}

\section{Interview questions}

\begin{interviewq}[$\star$ \iqroles{developer}]\label{iq:pl:testing-and-continuous-delivery:1}
What is the \omterm{def:pl:testing-and-continuous-delivery:pyramid}{test pyramid}, and why is it a pyramid?
\end{interviewq}

\begin{interviewq}[$\star\star$ \iqroles{developer}]\label{iq:pl:testing-and-continuous-delivery:2}
How do you regression-test a Monte Carlo pricer?
\end{interviewq}

\begin{interviewq}[$\star\star$ \iqroles{developer}]\label{iq:pl:testing-and-continuous-delivery:3}
A release reached seven of eight servers. How would your deployment process have stopped it?
\end{interviewq}

\begin{interviewq}[$\star\star$ \iqroles{developer}]\label{iq:pl:testing-and-continuous-delivery:4}
Would you deploy trading code continuously or on a \omterm{def:pl:testing-and-continuous-delivery:train}{release train}? Why?
\end{interviewq}

\begin{interviewq}[$\star\star\star$ \iqroles{developer}]\label{iq:pl:testing-and-continuous-delivery:5}
Design the test and release process of a firm's pricing library and order path.
\end{interviewq}
